\section{A learned aggregation of EPSS history}\label{app:learned}

\paragraph{Question.} LEV combines a CVE's EPSS history with one fixed rule.
Would a model that sees the same history, and nothing else, combine it better?
We restrict the inputs to EPSS history so that the comparison tests LEV's rule
rather than the value of additional data.

\paragraph{Design.} There are seven features, all computed from
version-normalized EPSS scores dated on or before the analysis date: LEV, log
CVE age, and the current, peak, mean, first and three-month-trend EPSS. There
are two learners: an L2-regularized logistic regression, with probability
features on the logit scale, and gradient-boosted trees, constrained to be
monotone in LEV and heavily regularized. The choice between them is fixed by a
rule that sees only training data: fit on the older two-thirds of the training
cohort, validate on the newer third, and keep the model with the higher average
precision (AP). We wrote this rule after a first, unreported run, and report
both models throughout. The baselines are LEV, current EPSS, peak EPSS, and the
isotonic and age-aware Platt maps of Section~\ref{sec:m-correct}. In Experiment~1, every
model, these maps included, is fitted on the same training set, the 10,296 CVEs
(10,255 once KEV-listed CVEs are removed) whose
first monthly EPSS snapshot is dated between 24 October 2024 and
30 January 2025; for the KEV target this
differs from Table~\ref{tab:prosp}, whose maps are fitted on the calibration
cohort, so their ratios differ too.

\emph{Experiment 1} repeats the prospective test of
Section~\ref{sec:m-correct} with three targets: the union of the primary lists;
CISA KEV; and the union with every KEV-listed CVE removed from both training
and test sets, which cuts the path by which a listing raises later scores. For
the logistic model, the bootstrap resamples training \emph{and} test CVEs and
refits, so its intervals include training variance.

\emph{Experiment 2} is a forward test. At an origin $t$ we take every CVE that
is at least 30 days old and on no list, compute its features at $t$, and label
it positive if it is first listed within the next 182 days. No feature reads
a snapshot dated after $t$. Training uses the
origin one horizon earlier. There are three test origins (2 August 2025,
31 October 2025 and 30 January 2026), and a second target that uses only CISA
and SENSOR dates, because VulnCheck's dates come from a later snapshot.

\begin{table}[t]
\centering\small
\caption{Experiment 1: average precision (and ratio of predicted to observed).
Fitted at 2025-08-01 and scored at 2026-08-01 on 44,316 unseen CVEs (44,124
once KEV-listed CVEs are removed), with 597, 192 and 405 test positives (111,
41 and 70 in training).}\label{tab:a1}
\begin{tabular}{@{}llll@{}}
\toprule
method & union & KEV & union without KEV \\
\midrule
LEV & 0.399 (2.03\x) & 0.324 (6.30\x) & 0.228 (2.66\x) \\
current EPSS & 0.404 (0.65\x) & 0.457 (2.02\x) & 0.221 (0.72\x) \\
peak EPSS & 0.424 (0.97\x) & 0.459 (3.03\x) & 0.237 (1.17\x) \\
isotonic & 0.337 (1.00\x) & 0.239 (1.29\x) & 0.203 (0.94\x) \\
Platt with age & 0.366 (0.88\x) & 0.329 (0.61\x) & 0.180 (1.20\x) \\
logistic & 0.398 (0.99\x) & 0.196 (1.52\x) & 0.190 (0.74\x) \\
boosted trees & 0.309 (1.05\x) & 0.127 (2.47\x) & 0.089 (0.45\x) \\
\bottomrule
\end{tabular}
\end{table}

\begin{figure}[tb]
\centering
\fig{learned_lev.pdf}
\caption{Left: Experiment 1 average precision by target; the learned models
never beat LEV or current EPSS. Right: Experiment 2 average precision by
forecast origin.}\label{fig:learned}
\end{figure}

\paragraph{No learned aggregation ranks better.} On every target, the logistic
model's AP is within noise of LEV's (Table~\ref{tab:a1}). The 95\% interval of
the difference is $-0.134$ to $+0.035$ on the union, $-0.137$ to $+0.021$ on
the union without KEV, and $-0.229$ to $+0.001$ on CISA KEV, where the logistic
model is also significantly below current EPSS ($-0.354$ to $-0.124$). The
boosted model is significantly worse than LEV on all three targets. Across
both experiments, the selection rule picked the logistic model in eight of nine
splits. Peak EPSS alone is as good as anything.

\paragraph{An apparent gain was an artifact.} A first version of this
experiment used EPSS scores that were not version-normalized and kept KEV-listed
CVEs in the target. There, the logistic model's AP exceeded LEV's (0.435
against 0.397), although the interval of the difference, once training
variance is included, spans zero ($-0.065$ to $+0.051$). The gap disappears
when either condition is removed. With version-normalized inputs, AP is 0.398
against 0.399. With KEV members dropped, it is 0.204 against 0.224 for LEV and
0.221 for current EPSS. A model trained on lists that feed EPSS learns the
feedback. The \texttt{E1\_raw\_epss} entry of \texttt{results/learned\_lev.json}
holds this check.

\begin{table}[t]
\centering\small
\caption{Experiment 2, union target: AP for first listing within the next 182
days, among unlisted CVEs at least 30 days old.}\label{tab:a2}
\begin{tabular}{@{}lrrrrrr@{}}
\toprule
test origin & positives / $N$ & base rate & LEV & current EPSS & logistic & boosted trees \\
\midrule
2025-08-02 & 71 / 32,223 & 0.22\% & 0.035 & 0.042 & 0.019 & 0.007 \\
2025-10-31 & 105 / 43,059 & 0.24\% & 0.028 & 0.027 & 0.038 & 0.025 \\
2026-01-30 & 97 / 49,678 & 0.20\% & 0.020 & 0.025 & 0.022 & 0.010 \\
\bottomrule
\end{tabular}
\end{table}

\paragraph{Nothing anticipates the lists.} Across the three forward origins
(Table~\ref{tab:a2}), no method's AP is significantly above that of current
EPSS. The boosted model is significantly below it at two origins and the
logistic model at one. Current EPSS is never significantly worse than LEV, and
is slightly better at two of the three origins when only CISA and SENSOR dates
are used. In each window, 59--75\% of the listings concern CVEs that were
less than 30 days old, or not yet scored, when the window opened, so no
history-based method can see them in advance.

\paragraph{Threats.} The training sets are small (20 to 111 positives),
Experiment 1 has a single split, and many intervals are reported without
correction for multiplicity, so these results are exploratory. SENSOR is an
EPSS label source, so some circularity remains even without KEV.

\section{Further robustness}\label{app:robust}

Each check below is regenerated by \texttt{make reproduce}, and its full
output is written to the file named.
\begin{itemize}
\item Reproduction error against analysis date (Section~\ref{sec:res-repro}):
\texttt{results/validate\_dn\_grid.json}.
\item Calibration by age and by maturity rule (Table~\ref{tab:age}):
\texttt{results/calibration.json}.
\item Growth families, HAC and block-bootstrap comparisons, and EPSS for
never-listed CVEs (Section~\ref{sec:res-growth}): \texttt{results/dn\_sweep.json}
and Figure~\ref{fig:sweepapp}.
\item The version artifact under percentile, step-only and cohort-fit maps,
and the one-day steps at v4 and v5 (Section~\ref{sec:res-version}):
\texttt{results/version\_normalize.json}.
\item The matched-control early-signal analysis (Section~\ref{sec:res-leak}):
\texttt{results/leakage.json}.
\end{itemize}

\begin{figure}[tb]
\centering
\fig{dn_sweep_appendix.pdf}
\caption{Analysis-date sweep: additional panels, and mean EPSS for CVEs that
no list records.}\label{fig:sweepapp}
\end{figure}
